%\documentclass[twocolumn]{article}\
\documentclass{article}
\usepackage{amsmath, amssymb, cancel, mathtools, bm}
\usepackage[left=.5in, right=.5in, top=1in, bottom=1in]{geometry}
\usepackage[most]{tcolorbox}
\usepackage[skip=1em,indent=0pt]{parskip}
\setlength{\parindent}{0pt}
\newcommand{\statvec}[1]{\underset{\sim}{\bm{#1}}} % Vector symbol (tilde under vector)
\DeclareMathOperator{\EX}{\mathbb{E}} % Expected Value symbol
\DeclareMathOperator{\ext}{\EX_{\theta}} % Expected Value symbol
\DeclareMathOperator{\vart}{Var_{\theta}} % Expected Value symbol
\newcommand{\indep}{\perp\!\!\!\!\perp} % Independence symbol
\newcommand{\real}{\mathbb{R}} % Simplified 'Reals' indicator
\newcommand{\pto}{\overset{P}{\to}}
\newcommand{\asto}{\overset{a.s.}{\to}}
\newcommand{\rthto}{\overset{L^r}{\to}}
\newcommand{\dto}{\overset{D}{\to}}
% Force all aggregate symbols to always put values above/below
\let\Oldint=\int
\let\Oldsum=\sum
\let\Oldprod=\prod
\let\Oldbigcup=\bigcup
\let\Oldbigcap=\bigcap
\let\Oldlim=\lim
\renewcommand{\int}{\Oldint\limits} 
\renewcommand{\sum}{\Oldsum\limits} 
\renewcommand{\prod}{\Oldprod\limits} 
\renewcommand{\bigcup}{\Oldbigcup\limits} 
\renewcommand{\bigcap}{\Oldbigcap\limits} 
\renewcommand{\lim}{\Oldlim\limits} 
\newcommand{\infint}{\int_{-\infty}^{\infty}}
\newcommand{\iiddist}{\overset{\mathrm{iid}}{\sim}}
\newcommand{\norm}[1]{\left\lVert#1\right\rVert}
\newcommand{\boldred}[1]{\textbf{\textcolor{red}{#1}}}
\usepackage[linktoc=all]{hyperref}



\begin{document}

\tableofcontents

\newpage

%%%%%%%%%%%%%%%%%%%%%%%%%%%%%%%%%%%%%%%%%%%%%%%%%%%%%%%%%%%%%%%%%%%%%%%%%%%
\section{Linear Algebra Review}

In stats: 

\begin{itemize}
    \item Matrices: the dataset
    \item Vector: one variable of the dataset (we will always assume vectors to be vertical)
    \item Scalar: single value of a variable
\end{itemize}

Matrix Multiplication:

$$
A = \begin{pmatrix} 1 & 2 & 3 \\ 4 & 5 & 6 \end{pmatrix}, 
B = \begin{pmatrix} 5 & 10 \\ 15 & 20 \\ 25 & 30 \end{pmatrix}, 
AB = \begin{pmatrix} (1*5 + 2*15 + 3*25) & (1*10+2*20+3*30) \\
    (4*5 + 5*15 + 6*25) & (4*10 + 5*20 + 6*30) 
\end{pmatrix} = 
\begin{pmatrix}
    110 & 140 \\ 245  & 320
\end{pmatrix}
$$

$AB \neq BA$

Identity Matrix is just a diagonla matrix of all 1's

Basic Operations:
\begin{itemize}
    \item Only square matrices have inverses
    \item Find $A^{-1}$ is long and slow
    \item $(A^{-1})^{-1} = A$
    \item $(sA)^{-1} = s^{-1}A^{-1}$ ($s$ is a scalar)
    \item $(AB)^{-1} = B^{-1}A^{-1}$
    \item If $D$ is diagonal, $D^{-1}$ is also diagonal with elements $\frac{1}{d_{ii}}$
\end{itemize}


The Determinant ($|A| = s$) corresponds to Variance of a dataset

Gradients:

$f(x_1, x_2, x_3) = 3x_1^2x_2 + x_2x_3^3 - e^{x_1x_2}, \nabla f(x) = 
\begin{pmatrix}
    6x_1x_2-x_3e^{x_1x_3} \\
    3x_1^2 + x_3^3 \\
    3x_2x_3^2 - x_1e^{x_1x_3}
\end{pmatrix}
$

\end{document}